%%%%%%%%%%%%%%%%%%%%%%%%
%
% $Autor: Hemanth Jadiswami Prabhakaran $
% $Datum: 2026-09-25 $
% $Pfad: GitHub/MBIDA_Master_Thesis/Manual/Chapters/en/07Maintenance.tex $
% $Version: 1 $
%
% $Project: MBIDA Thesis - Manual $
%
%%%%%%%%%%%%%%%%%%%%%%%%


\chapter{Maintenance}
\label{chap:maintenance}

This chapter is written for the \emph{operator}: the person who owns the GitHub repository and the Streamlit Community Cloud app. Users of the dashboard do not need it.

Because the data is static, maintenance is small. It consists of keeping the online app available, keeping the software working, and redeploying after a change. The forecasting pipeline is \emph{not} part of regular maintenance; it is only run again if the result files are lost (Section~\ref{sec:mtRegenerate}).

\section{Maintenance Cycle}
\label{sec:mtCycle}

Every change to the dashboard -- a new package version, a corrected text, a changed setting -- follows the same cycle (Figure~\ref{fig:redeployCycle}). The online app is rebuilt automatically whenever the repository's main branch changes; there is no separate upload step.

\begin{figure}[H]
  \centering
  \resizebox{0.95\linewidth}{!}{% Maintenance cycle for the dashboard (app only - the data is static).
\begin{tikzpicture}[
    st/.style={man/box,text width=3.0cm,minimum height=1.25cm,font=\sffamily\small},
  ]
  \def\R{3.6}
  \node[st]                   (a) at (90:\R)  {\textbf{1} \; Change\\\footnotesize code, package pin, config};
  \node[st]                   (b) at (18:\R)  {\textbf{2} \; Test locally\\\footnotesize\texttt{streamlit run}};
  \node[st]                   (c) at (-54:\R) {\textbf{3} \; Commit \& push\\\footnotesize to \texttt{main}};
  \node[st,man/boxgreen]      (d) at (-126:\R){\textbf{4} \; Cloud rebuilds\\\footnotesize automatically};
  \node[st]                   (e) at (162:\R) {\textbf{5} \; Verify online\\\footnotesize checklist, Sec.~\ref{sec:mtVerify}};

  \draw[man/arrow] (a) to[bend left=18] (b);
  \draw[man/arrow] (b) to[bend left=18] (c);
  \draw[man/arrow] (c) to[bend left=18] (d);
  \draw[man/arrow] (d) to[bend left=18] (e);
  \draw[man/arrow] (e) to[bend left=18] (a);

  \node[man/terminal,text width=2.3cm,font=\sffamily\small] (mid) at (0,0) {Dashboard\\in service};

  % exception branch: result files missing
  \node[man/boxred,text width=3.4cm,font=\sffamily\small] (x) at (7.9,-0.2)
     {\textbf{Only if the 4 result files are missing:}\\\footnotesize run the pipeline notebook once (Sec.~\ref{sec:mtRegenerate})};
  \draw[man/dashedarrow] (b.east) to[out=0,in=90] (x.north);
  \draw[man/dashedarrow] (x.south) to[out=-90,in=0] (c.east);
\end{tikzpicture}
}
  \caption{Maintenance and redeployment cycle}
  \label{fig:redeployCycle}
\end{figure}

\begin{enumerate}
  \item \textbf{Change} the file in a local copy of the repository.
  \item \textbf{Test locally} with \texttt{streamlit run Code/app/streamlit\_app.py} (Section~\ref{sec:instLocal}) and do the verification in Section~\ref{sec:mtVerify}.
  \item \textbf{Commit and push} to the main branch:
\begin{lstlisting}[style=shell]
git add <changed files>
git commit -m "Short description of the change"
git push
\end{lstlisting}
  \item \textbf{Wait for the rebuild.} Streamlit Community Cloud notices the push and restarts the app within a few minutes; if packages changed, it installs them again first.
  \item \textbf{Verify online} with the same checklist.
\end{enumerate}

\section{Routine Tasks}
\label{sec:mtRoutine}

\begin{table}[H]
\centering
\caption{Routine maintenance tasks}
\label{tab:mtRoutine}
\small
\begin{tabularx}{\linewidth}{@{}p{3.1cm}Xp{3.0cm}@{}}
\toprule
\textbf{When} & \textbf{Task} & \textbf{Reference} \\
\midrule
Before a presentation or an exam & Open the online app some minutes in advance so that it is awake; run the verification checklist & Sections~\ref{sec:instWake}, \ref{sec:mtVerify} \\
Monthly & Open the app once; check the Streamlit Community Cloud workspace for warnings or e-mails & Section~\ref{sec:mtCloud} \\
Every 6 months & Check whether the pinned package versions still install; update if needed & Section~\ref{sec:mtDeps} \\
After every change & Complete maintenance cycle & Section~\ref{sec:mtCycle} \\
If files are lost & Regenerate the result files & Section~\ref{sec:mtRegenerate} \\
\bottomrule
\end{tabularx}
\end{table}

\section{Managing the Online App}
\label{sec:mtCloud}

The online app is managed in the Streamlit Community Cloud workspace at \url{https://share.streamlit.io}, after signing in with the GitHub account that owns the repository \cite{StreamlitCloud:2026}. The workspace lists the app with its status. The menu next to it, and the \UI{Manage app} button at the bottom right of the running app (visible only to the signed-in owner), offer:

\begin{table}[H]
\centering
\caption{Actions in the Streamlit Community Cloud workspace}
\label{tab:mtCloudActions}
\small
\begin{tabularx}{\linewidth}{@{}lX@{}}
\toprule
\textbf{Action} & \textbf{Use} \\
\midrule
Logs & Show the app's console output, e.g.\ package installation errors or Python error messages \\
Reboot app & Restart the app without changing anything -- first aid for a hanging app \\
Settings & Change the app's address and sharing settings \\
Delete & Remove the app (the repository is not affected) \\
\bottomrule
\end{tabularx}
\end{table}

\begin{figure}[H]
  \centering
  \Screenshot[0.85\linewidth]{Images/07Maintenance/CloudManageApp.png}
  \caption{The \UI{Manage app} panel with logs and app menu}
  \label{fig:mtCloudManage}
\end{figure}

The deployment itself is configured as follows. These values are set once when the app is created; they only need to be checked if the app is ever deployed anew.

\begin{table}[H]
\centering
\caption{Deployment settings of the online app}
\label{tab:mtDeploySettings}
\small
\begin{tabularx}{\linewidth}{@{}lX@{}}
\toprule
\textbf{Setting} & \textbf{Value} \\
\midrule
Repository & \texttt{Hemanth-Jp/MBIDA\_Master\_Thesis} \\
Branch & \texttt{main} \\
Main file path & \texttt{Code/app/streamlit\_app.py} \\
App \ac{url} & \texttt{m5-hierarchical-forecast-explorer} \\
Python version & 3.10 or newer (chosen under \UI{Advanced settings}) \\
Package list & \texttt{Code/app/requirements.txt} (found automatically next to the main file) \\
\bottomrule
\end{tabularx}
\end{table}

\begin{WarningBox}
The Python version can only be chosen when the app is created. To change it, delete the app and deploy it again with the same settings and the same \ac{url}.
\end{WarningBox}

\section{Updating Packages}
\label{sec:mtDeps}

The five packages are pinned to exact versions in \texttt{Code/app/requirements.txt}. Pinning guarantees that the online app and every local installation behave identically. To update:

\begin{enumerate}
  \item In a local copy, change one version number at a time, e.g.\ \texttt{streamlit==1.61.1} to a newer release.
  \item Reinstall: \texttt{pip install -r Code/app/requirements.txt}.
  \item Start the app and run the verification checklist (Section~\ref{sec:mtVerify}), including light and dark mode.
  \item If everything works, complete the maintenance cycle; otherwise return to the previous version number.
\end{enumerate}

\begin{WarningBox}
Keep the package name \texttt{st-theme} exactly as written. The similarly named \texttt{streamlit-theme} is a different package, and the dashboard fails to start with it.
\end{WarningBox}

\section{Changing Settings}
\label{sec:mtConfig}

Names shown in the sidebar -- the list of base models and the reconciliation methods with their labels -- come from the settings file \texttt{Code/config.py}, which the dashboard shares with the forecasting pipeline. Two rules apply:
\begin{itemize}
  \item Visible \emph{labels} of reconciliation methods (for example \VAL{MinTrace (wls\_struct)}) may be renamed freely.
  \item The technical \emph{names} behind them must match the column names in the result files exactly. If they do not, the dashboard reports ``No column found'' (Section~\ref{sec:tsNoColumn}). They only need to change if the pipeline itself changes.
\end{itemize}

The look of the Streamlit page (colours, fonts) is set in \texttt{.streamlit/config.toml} in the repository root.

\section{Verification Checklist}
\label{sec:mtVerify}

Run this list after every change, locally before pushing and online after the rebuild.

\begin{todolist}
  \item The page loads without a red error box.
  \item The node line shows \VAL{Total (level = 1)} at start.
  \item The chart shows the actual line, a dotted base line and a solid \VAL{MinTrace} line, and the three shaded regions.
  \item Drilling down to \VAL{CA} $\rightarrow$ \VAL{CA\_1} $\rightarrow$ \VAL{FOODS} $\rightarrow$ \VAL{FOODS\_3} $\rightarrow$ \VAL{FOODS\_3\_090} works; the level increases to 6.
  \item Each of the three base models changes the dotted line.
  \item All three reconciliation methods can be shown together; no ``No column found'' text appears.
  \item Light and dark mode both give a readable chart.
  \item The camera button saves a \ac{png}.
\end{todolist}

\section{Regenerating the Result Files}
\label{sec:mtRegenerate}

This is needed only if the four files in \texttt{Code/artifacts/dashboard\_data/} are missing or damaged, which the dashboard reports with a red error box (Section~\ref{sec:tsParquet}). First try to restore them from Git, which is much faster:

\begin{lstlisting}[style=shell]
git checkout -- Code/artifacts/dashboard_data/
\end{lstlisting}

If that is not possible, recreate them with the thesis pipeline. Because the complete daily data set is processed, this can take half an hour or more and needs a computer with plenty of memory (16\,GB of \ac{ram} recommended).

\begin{enumerate}
  \item Download \texttt{sales\_train\_validation.csv} and \texttt{calendar.csv} from the \ac{m5} competition page \cite{Kaggle:M5} into \texttt{Code/data/m5/} (see \texttt{Code/data/README.md}).
  \item Install the pipeline packages in a separate environment:
\begin{lstlisting}[style=shell]
python3 -m venv .venv-pipeline
source .venv-pipeline/bin/activate      # Windows: .venv-pipeline\Scripts\Activate.ps1
pip install -r Code/requirements-pipeline.txt
\end{lstlisting}
  \item Open \texttt{Code/notebooks/02\_pipeline.ipynb} in Jupyter and run \emph{all cells from top to bottom} (\UI{Run} $\rightarrow$ \UI{Run All Cells}). The last cell writes the four files.
  \item Check that they exist, then commit and push them (Section~\ref{sec:mtCycle}):
\begin{lstlisting}[style=shell]
ls Code/artifacts/dashboard_data
git add Code/artifacts/dashboard_data
git commit -m "Regenerate dashboard result files"
git push
\end{lstlisting}
\end{enumerate}

\begin{NoteBox}
With unchanged input data and settings, the pipeline reproduces the same results. The online app only shows new files after they have been pushed -- running the pipeline locally alone does not change the online dashboard.
\end{NoteBox}

\section{Backup and Rollback}
\label{sec:mtBackup}

The Git history is the backup: every version of the program and of the result files is stored in the repository. To undo a change that broke the online app:

\begin{lstlisting}[style=shell]
git log --oneline            # find the faulty commit
git revert <commit-id>       # create a commit that undoes it
git push                     # the online app rebuilds with the old state
\end{lstlisting}

\section{Resetting a Local Installation}
\label{sec:mtLocalReset}

If a local installation behaves strangely, rebuilding the virtual environment is quick and safe; it does not touch the repository files:

\begin{lstlisting}[style=shell]
deactivate                    # if an environment is active
rm -rf .venv                  # Windows: Remove-Item -Recurse -Force .venv
python3 -m venv .venv
source .venv/bin/activate
pip install -r Code/app/requirements.txt
\end{lstlisting}
